\chapter{Penutup}
\section{Kesimpulan}
Dalam pengerjaan tugas akhir ini, telah ditinjau performa 3 siklus mesin panas kuantum fraksional; Carnot, Otto, dan Brayton. Peninjauan dilakukan untuk efisiensi dan juga \textit{power output} terhadap parameter ekspansi, fraksional, dan disipasi. Variasi dilakukan terhadap parameter fraksional dan disipasi di mana ketika parameter fraksional berubah, parameter  disipasi dibuat konstan dan berlaku sebaliknya. Dari peninjauan yang dilakukan, didapatkan kesimpulan sebagai berikut:
\begin{enumerate}
    \item Parameter fraksional dapat diimplementasi ke permasalahan mesin panas kuantum dengan mengganti persamaan Schrodinger yang digunakan menjadi persamaan Schrodinger fraksional. Persamaan Schrodinger fraksional ini kemudian digunakan untuk mendapatkan besaran energi eigen sesuai dengan syarat batas yang berlaku. Pada peninjauan ini adalah permasalahan potensial tak hingga.
    \item Parameter fraksional berpengaruh terhadap efisiensi serta \textit{power output} dari setiap siklus di mana penurunan nilai \(\alpha\) akan berakibat penurunan efisiensi serta \textit{power output}. Hanya saja pengaruh \(\alpha\) lebih signifikan terhadap \textit{power output} dibanding terhadap efisiensi.
    \item Parameter disipasi berpengaruh cukup signifikan terhadap efisiensi. Pada keadaan tanpa disipasi, dimungkinkan untuk mencapai efisiensi 100\%. Pada keadaan dengan disipasi, hal tersebut tidak dimungkinkan lagi. Lalu, kenaikan nilai  \(\beta\) akan berakibat penurunan efisiensi secara signifikan. Meski demikian, \(\beta\) tidak memiliki pengaruh apapun terhadap \textit{power output} sehubuah siklus mesin panas kuantum.
\end{enumerate}

\section{Saran}
Pada penelitian ini dilakukan inkorporasi kalkulus fraksional terhadap partikel non-relativistik. Peninjauan dapat dilakukan untuk partikel relativistik. Selain itu, peninjauan dilakukan untuk partikel 1D yang terjebak pada potensial tak hingga. Penelitian berikutnya dapat juga dilakukan untuk mengekspasi dimensi yang ditinjau hingga 3D. Terakhir, siklus yang ditinjau pada tulisan ini adalah Carnot, Otto, dan Brayton. Tinjauan untuk siklus-siklus mesin panas kuantum fraksional lain juga dapat disarankan untuk dilakukan. 

\newpage
\thispagestyle{empty}
\vspace*{\fill}
    \begin{center}
	Halaman ini sengaja dikosongkan
    \end{center}
\vspace*{\fill}
\pagebreak